\documentclass{article}


\usepackage[a4paper, margin=2cm]{geometry}
%
\usepackage{amsfonts, amsmath, amssymb}
\usepackage{mathrsfs, mathtools, mathbbol}
\usepackage{tikz}
\usepackage{tikz-cd}
%
\usepackage[T2A]{fontenc}
\usepackage[utf8]{inputenc}
\usepackage[english]{babel}
\usepackage{alphabeta}
\usepackage{hyperref}

\usepackage[backend=biber]{biblatex}
\usetikzlibrary{shapes,arrows, positioning,fit}

\newcommand*{\sheaf}{\mathscr}


\usepackage[most]{tcolorbox}

\newtcbtheorem[number within=section]{definition}{Definition}{
    enhanced,
    breakable,
    colback=gray!5,
    colframe=gray!50,
    fonttitle=\bfseries,
    separator sign={.},
    before skip=10pt,
    after skip=10pt,
    boxrule=0.5pt,
    arc=2pt,
    left=8pt,
    right=8pt,
    top=6pt,
    bottom=6pt
}{def}


\newtcbtheorem[number within=section]{proposition}{Proposition}{
    enhanced,
    breakable,
    colback=gray!5,
    colframe=gray!50,
    fonttitle=\bfseries,
    separator sign={.},
    before skip=10pt,
    after skip=10pt,
    boxrule=0.5pt,
    arc=2pt,
    left=8pt,
    right=8pt,
    top=6pt,
    bottom=6pt
}{prop}


\newtcbtheorem[number within=section]{example}{Example}{
    enhanced,
    breakable,
    colback=gray!5,
    colframe=gray!50,
    fonttitle=\bfseries,
    separator sign={.},
    before skip=10pt,
    after skip=10pt,
    boxrule=0.5pt,
    arc=2pt,
    left=8pt,
    right=8pt,
    top=6pt,
    bottom=6pt
}{ex}


\newtcbtheorem[number within=section]{remark}{Remark}{
    enhanced,
    breakable,
    colback=gray!5,
    colframe=gray!50,
    fonttitle=\bfseries,
    separator sign={.},
    before skip=10pt,
    after skip=10pt,
    boxrule=0.5pt,
    arc=2pt,
    left=8pt,
    right=8pt,
    top=6pt,
    bottom=6pt
}{remark}


\newtcbtheorem[number within=section]{note}{Note}{
    enhanced,
    breakable,
    colback=gray!5,
    colframe=gray!50,
    fonttitle=\bfseries,
    separator sign={.},
    before skip=10pt,
    after skip=10pt,
    boxrule=0.5pt,
    arc=2pt,
    left=8pt,
    right=8pt,
    top=6pt,
    bottom=6pt
}{note}


\newtcbtheorem[number within=section]{conjecture}{Conjecture}{
    enhanced,
    breakable,
    colback=gray!5,
    colframe=gray!50,
    fonttitle=\bfseries,
    separator sign={.},
    before skip=10pt,
    after skip=10pt,
    boxrule=0.5pt,
    arc=2pt,
    left=8pt,
    right=8pt,
    top=6pt,
    bottom=6pt
}{conjecture}



\newtcbtheorem[number within=section]{question}{Question}{
    enhanced,
    breakable,
    colback=gray!5,
    colframe=gray!50,
    fonttitle=\bfseries,
    separator sign={.},
    before skip=10pt,
    after skip=10pt,
    boxrule=0.5pt,
    arc=2pt,
    left=8pt,
    right=8pt,
    top=6pt,
    bottom=6pt
}{question}






\addbibresource{ainfty-geometry.bib}

\title{$A_\infty$ Enhancement of Geometry}
\author{Koji Sunami}

\begin{document}


\maketitle
\vspace{1cm}


The goal of our research is to prove homological mirror symmetry conjecture or to develop the relevant area of study.
The claim of homological mirror symmetry conjecture is $A_\infty$-equivalence of Fukaya category $Fuk(X)$ 
and $A_\infty$-enhancement $D^b(Coh(X))$ (typically noted as $D^b(X)$) of derived category of coherent sheaves.
This statement says the elaboration of mirror symmetry, originally correspondence of symplectic data on Hodge structure and that of complex
in Calabi-Yau manifold, whose symplectic side we name as A-model, the complex side as B-model, 
mirror symmetry solves many of QFT problem as well as Landau-Ginzburg model and others,
and homological mirror symmetry is generalization of mirror symmetry in $\infty$-category. 
The proof isn't easy, but the understanding of the problem statement itself is not difficult. \\


\begin{definition}{$A_\infty$ Category}{}
An $A_\infty$ category $\mathcal{A}$ is a quasi-category which consists of objects $X\in Ob(\mathcal{A})$
and $1$-morphisms $f\in Mor^1(\mathcal{A})$, and the $1$-morphism contains higher coherence $\alpha: f\Rightarrow g$, 
called 2-morphism, and iteratively $n$-morphisms.
\end{definition}


\begin{definition}{$A_\infty$ Functor}{}


\end{definition}


\begin{definition}{Fukaya Category}{}
For a symplectic manifold $X$, Fukaya category $Fuk(X)$ is an $A_\infty$ category describing Lagrangian intersection.
The objects of Fukaya category $Ob(Fuk(X))$ is Lagrangian submanifolds $L\subset X$, 
and the morphisms are Floer chain complex $CF^*(L_0,L_1)$ for each $L_1$ and $L_2$.
This $CF^*(L_0,L_1)$ is a $\mathbb{Z}$-graded vector space defined by $CF^*(L_0,L_1) = \underset{p\in L_0\cap L_1}{\oplus} \mathbb{k} p$.
Floer chain complex is generated by its intersection points, and this is equivalent with counting critical points of Morse theory. \\

In Fukaya category, there is a composition $m_k: CF^*(L_{k-1},L_k)\otimes \cdots \otimes CF^*(L_1,L_2) \to CF^*(L_0,L_k)[2-k]$
satisfying Stasheff associahedron property.
\end{definition}


$A_\infty$ enhancement of dg-category doesn't have unique enhancement, 
but the uniqueness is guaranteed if dg-category is derived category $D^b(X)$.

\begin{definition}{$A_\infty$ Derived Category}{}
$A_\infty$ enhancement of derived categeory $D^b(X)$ is an $A_\infty$ category $\mathcal{A}$,
whose homology is $H^*(\mathcal{A}(E,F))=Ext^*(E,F)$, 
while its trivialization $H^0(\mathcal{A}(E,F))$ needs to become the based derived category $H^0(\mathcal{A}(E,F)) \cong Hom_{D^b(X)}(E,F)$.
This $A_\infty$ enhancement $\mathcal{A}$ can be denoted as $D^b(X)$ by abuse of notation.
\end{definition}


\begin{conjecture}{Homological Mirror Symmetry}{}
There is an $A_\infty$ equivalence of categories 
\[
Fuk(X) \cong D^b(X)
\]
\end{conjecture}




What is the structure of $A_\infty$ category? The natural question is geometric realization of $A_\infty$ category.
Starting from $A_\infty$ algebra as a special case of $A_\infty$ category, I define it as noncommutative formal pointed dg-manifold
(hereafter in short, ncfp dg-manifold), and then I'll generalize to $A_\infty$ category.
Briefly the basic idea is $A_\infty$ algebra derives coalgebra, and coalgebra is contravariantly mapped to algebra in category,
and this is an $A_\infty$ extnesion of notion of affine scheme.
In fact, the case of $A_\infty$ category is analogous to $A_\infty$ algebra.


\begin{definition}{Counital Coalgebra}{counital coalgebra}
Coalgebra $B$ is a graded vector space with morphism $\Delta:B\to B\otimes B$ s.t. 
$(\Delta \otimes id)\circ \Delta = (id \otimes \Delta)\circ \Delta$ called coproduct,
and counital coalgebra is coalgebra $B$ with unit $\epsilon:B\to 1$ where $1$ is the unit object of $\mathcal{C}$,
satisfying $(\epsilon \otimes id)\circ \Delta = (id \otimes \epsilon)\circ \Delta$.
Counital coalgebra $B$ is cofree counital coalgebra if it's given by $B=C(V)\cong T^c(V)$.
\end{definition}


\begin{definition}{A Non-Commutative Thin Scheme}{a noncommutative thin scheme}
A non-commutative thin scheme $X$ is a functor $X:Alg_{\mathcal{C}^f} \to Sets$, that is isomorphic to $X\cong Hom(-^*, B_X)$
for some coalgebra $B_X$, where $\mathcal{C}^f$ is artinian category of finite dimensional objects in $\mathcal{C}$, 
where $\mathcal{C}$ is $\mathbb{Z}$-graded vector spaces.
In particular, $X$ is a non-commutative formal manifold, if the functor $X$ is isomorphic to some $Spc(T(V))$
for a cofree counital tensor coalgebra $T(V)=\underset{\geqq 1}{\oplus}V^{\otimes n}$. The dimension of $X$ is defined as $dim_k V$.
\end{definition}


\begin{definition}{A Non-Commutative Thin DG-Scheme}{} 
A non-commutative thin dg-scheme is a pair $(X,d)$ where $X$ is a non-commutative thin scheme, and $d$ is a vector field on $X$
of degree $+1$ satisfying bracket of dgla as $[d,d]=0$. This $d$ is called homological vector field.
\end{definition}


\begin{definition}{Non-Commutative Formal Pointed DG-Manifold}{}
A non-commutative formal pointed dg-manifold is a pair $((X, x_0), d_X)$ such that $(X, x_0)$ 
is a non-commutative formal pointed graded manifold, and $d_X$ is a homological vector field on $X$ such that $d_X|_{x_0} = 0$. 
\end{definition}


\begin{proposition}{Taylor Decomposition of Graded Vector Field}{}
For the case of formal manifold $X$, let the corresponding cofree counital coalgebra be $B_X\cong T^c(V)$,
and the homological vector field be coderivation $d_X:B_X\to B_X$, that can have restriction by cofreeness, thus by using projection,
$\pi \circ d_X:B_X\to V[1]$. If in addition, $d_X$ can be decomposed degreewise for its linearity,
as $d_X = \underset{n\geqq 1}{\Sigma} d_X^{(n)}$ for linear maps $d_X^{(n)}:V^{\otimes n}\to V[1]$.
\end{proposition}


\begin{definition}{Non-Unital $A_\infty$ Algebra}{}
A non-unital $A_\infty$-algebra $A=(A,m_n)$ is a graded vector space $A$ and coherence condition $m_n: A^{\otimes n}\to A[2-n]$ 
with Stasheff relation, whose geometric interpretation is the structure of a nc fp dg-manifold $((X, x_0), d_X)$ 
together with a coalgebra isomorphism $B_X\cong T(T_{x_0}X)$ with $x_0:k\to B_X$.
\end{definition}


\begin{definition}{Structure of $A_\infty$ Algebra}{}
The structure of $A_\infty$ Algebra on $V\in Ob(Vect_k^{\mathbb{Z}})$ is given by coderivation $d$ of degree $+1$ 
of the non-counital cofree coalgebra $T_{+}(V[1]) = \underset{n\geqq 1}{\oplus}V^{\otimes n}$ s.t. $[d,d]=0$.
\end{definition}


\begin{proposition}{Stasheff Relation}{}
The condition of nc fp dg-manifold $((X,x_0),d_X)$ derives Stasheff relation of the $A_\infty$ algebra,
that is $m_n:= d_X^{(n)}$ for $m_n: A^{\otimes n}\to A$ s.t.
$\underset{r+s+t=n}{\Sigma} (-1)^{r+st} m_{r+1+t} (id^{\otimes r} \otimes m_s \otimes id^{\otimes t})= 0$, whose degree is  $|m_n|=2-n$.
The Taylor decomposition of homological vector field $\pi\circ d_X = \underset{n\geqq 1}{\Sigma} d_X^{(n)}$ is degreewise by linearity,
as $d_X^{(n)}:B_X^{(n)} \to B_X^{(1)}$.
\end{proposition}


\begin{definition}{Non-Commutative Formal n-Pointed DG-Manifold}{}
Nc formal $n$-pointed dg-manifold $(\mathcal{X}, d_{\mathcal{X}}, \{pt_i\})$ is nc formal graded dg-manifold $(\mathcal{X},d_{\mathcal{X}})$,
and the points are $pt_i \in \mathcal{X} (k)$ for $1 \leqq i \leqq n$, and coderivation satisfies $d_{\mathcal{X}}|_{pt_i}=0$,
$d_{\mathcal{X}}|_{TX_i}: TX_i \to  \underset{j,k}{\oplus} TX_j \oplus TX_k$ where $X_i$ is formal neighborhood of $X$ at $pt_i$.
\end{definition}


\begin{definition}{$A_\infty$-Category}{}
$A_\infty$-category $\mathcal{A}$ is nc formal dg-manifold $(\mathcal{X}, d_{\mathcal{X}}, \{pt_i\})$
with its coalgebra $B_{\mathcal{X}}\cong T^c (V[1])$. This $V$ has a decomposition $V= \underset{i,j}{\oplus} E_{ij}$,
where $E_{ij}:= T(\widehat{X_i \times X_j}_{pt_i,pt_j})$.
\end{definition} 


\begin{remark} [] 
This $V$ is a $k$-linear associative algebra generated from $k$-linear category $\mathcal{C}$: 
let $Ob(\mathcal{C})=I$ to define an algebra $V= \oplus Hom(i,j)$ for $i,j\in I$ is a $k$-linear algebra
with the unit $1_V = \underset{i\in I}{\Sigma} id_i$, where $\pi_i = id_i\in V$ is a commuting idempotent i.e. $\pi_i^2 = \pi_i$, 
$\pi_i \pi_j = \pi_j \pi_i$. This makes homomorphism $k^I \to V$ of algebra with unit.
\end{remark}


\begin{proposition}{Stasheff Relation}{}
From non-unital $A_\infty$ category $\mathcal{A}$, the coalgebra $d_{\mathcal{X}}: T^c (\oplus E_{ij}[1]) \to T^c (\oplus E_{ij}[1])$ 
satisfies the Stasheff relation. Consider the tensor coalgebra can be given by the grading structure
$T^c(\underset{1\leqq i,j\leqq n}{\oplus} E_{ij}[1]) = \underset{m\geqq 0}{\oplus} \underset{1\leqq i_0,i_1,i_2,\cdots i_m\leqq n}{\oplus}
E_{i_0i_1} \otimes E_{i_1i_2} \otimes \cdots \otimes E_{i_{n-1}i_n} := \underset{m\geqq 0}{\oplus} B_{\mathcal{X}}^m$.
From cofreeness of the coalgebra, the coderivation $d_{\mathcal{X}}$ canonically projected onto $\pi:B_{\mathcal{X}} \to B_{\mathcal{X}}^1$,
and that can be decomposed linearwise, so $\pi \circ d_{\mathcal{X}} = \Sigma d_{\mathcal{X}}^n$ where
$d_{\mathcal{X}}^n: B_{\mathcal{X}}^n \to B_{\mathcal{X}}^1$, and we have Stasheff relation from $[d_{\mathcal{X}}, d_{\mathcal{X}}]=0$,
and $m_n := d_{\mathcal{X}}^n$.
\end{proposition}


\begin{remark} {}{}
$d_{\mathcal{X}}\circ \eta = 0$ implies $d_{\mathcal{X}}\circ \eta (1) = 0$ or $d_{\mathcal{X}} (1) = 0$. This means curvature is $m_0=0$.
\end{remark}


\begin{remark}{}{}
For the case of $I=\{i\}$ consists of one vertex, we write $d|_{pt} = d\circ \eta$.
\end{remark}




%\mathcal{}}$ with isomorphism $B_{\mathcal{X}}\cong T^c(Q[1])$ for a quiver
%$Q=\underset{X,Y\in Ob(\mathcal{A})}{\oplus} Hom(X,Y)$ is a quiver, $T^c(Q[1])=\underset{n\geqq 0}{\oplus} Q[1]^{\otimes n}$
%is cofree coassociative coalgebra, and source/target needs to be corresponding in this grading.
%This $[1]:Q\to Q[1]=\underset{X,Y}{\oplus}Hom(X,Y)[1]$ is a morphism of shifting each direct summand $Hom(X,Y)$, which is a graded vector space.
%The homological vector field $d_{\mathcal{X}}$ is a coderivation $d_{\mathcal{X}}: T^c(Q[1])\to T^c(Q[1])$ s.t. 
%$[d_{\mathcal{X}}, d_{\mathcal{X}}] = 0$.
%In particular if $\mathcal{A}$ consists one object $X$, $Q=Hom(X,X)$ is canonically non-unital $A_\infty$-algebra.
%\end{definition}



%\begin{definition} (The Structure of $U_f$) \\
%To describe its structure of nc formal neighborhood $U_f = \hat{Map(X,Y)}_f$, this is nc formal pointed manifold pointed at $f\in Maps(X,Y)$,
%whose corresponding $A_\infty$ algebra is $Center(f)=\underset{n\geqq 1}{\oplus} Hom(B_X^{\otimes n},B_Y)[n]$ as a graded vector space,
%where $B_X$ and $B_Y$ are the corresponding $A_\infty$ algebra, 
%and its Stasheff relation $(m_n)$ is derived from the homological vector field $d_Z$ where $Z=Hom(X,Y)$ s.t. $d_Z|_f = 0$,
%and this gives rises to $d_Z = d_X \otimes 1_Y - 1_X \otimes d_Y$. To describe Stasheff relation $(m_n)$ from $d_Z$, 
%$m_n(\phi_1,\cdots, \phi_n)(a_1,a_2,\cdots, a_N)=I+R$ for $\phi_1,\phi_2,\cdots, \phi_n\in Center(f)$ and $a_1,a_2,\cdots, a_N\in A$.
%$I$ and $R$ are compositions of different tree diagrams to assign the partititon of $\phi_i$ and $a_j$ to define $m_n$,
%and $I$ corresponds to the term $1_X\otimes d_Y$ in other words $d_Y\circ f$, and $R$ corresponds to the term $d_X\otimes 1_Y$.
%Partitioning $1\leqq i_1 < i_2 < \cdots < i_n = N$, and $b_k = \phi_k(a_{i_{k-1}+1},\cdots, a_{i_k})\in B_Y$, 
%whose composition is $m_n^B (b_1,\cdots, b_n)$, while the other partitioning $1\leqq j_1 < j_2 < \cdots < j_n = N$
%defines $c_k =m_{j_k-j_{k-1}}^A(a_{j_{k-1}+1},\cdots, a_{j_k})\in B_X$, whose composition is $\phi_1(c_1,\cdots, c_n)\in B_Y$,
%but $R$ appears only for $n=1$. Adding up the whole possibility of partitioning defines $I$ and $R$ respectively to define $m_n$.
%\end{definition}
%
%
%\begin{proposition} (Stasheff Relation) \\
%$m_n$ satisfies Stasheff relation.
%$[d_{\mathcal{X}}, d_{\mathcal{X}}] = 0$ makes $[\pi\circ d_{\mathcal{X}}, \pi\circ d_{\mathcal{X}}] = 0$,
%equivalently $[\underset{n\geqq 1}{\Sigma} m_n, \underset{n\geqq 1}{\Sigma} m_n]=0$.
%The proof is the same as $A_\infty$ algebra case.
%\end{proposition}
%
%For the case of infinite objects $A_\infty$ category is its inductive limit of finite objects.


\vspace{0.5cm}
\noindent\rule{\textwidth}{0.4pt}
\vspace{0.5cm}

I'll define non-commtuative scheme from its coalgebra identification of $A_\infty$ algebra/category.
Consider a functor $X: Alg \to Coalg$, mapping an associative algebra $O(X)$ to a coalgebra $B_X$.
Associative algebra is more general than unital commutative associative algebra $CRing$, 
and it's different from classical commutative scheme. Some idea of understanding noncommutative scheme is \\


Recall a commutative ring defines an affine scheme by its ring spectrum $X=V(\frak{p})$ with its structure sheaf $\mathcal{O}_X$, 
and that can also be described functorially as $F: CRing \to Set$. Hereafter I focus on functorial point of view. \\


\begin{itemize}
\item Deformation
\item Quantization
\end{itemize}

I don't know the general classification of associative algebra,
but one important fact is that commutative algebra's quantization is associative.
I wonder if there is any connection to the classical algebraic geometry or derived algebraic geometry in general. \\

\begin{remark}{Associative Algebra Example}{}
To name some examples of associative algebra, $U(\frak{g})$, $k[G]$, $M_n(k)$, Hecke algebra, path algebra, quiver algebra etc,
and quantization of associtive algebra is associtive algebra.
\end{remark}

\begin{question}{}{}
Should we just use deformation of associative algebra, which makes $A_\infty$ algebra?
\end{question}


Another approach is deformation.
One idea for the further development is arguing the correspondece of $A_\infty$ algebra and $L_\infty$ algebra through Hochschild cohomology 
through Kontsevich-Soibelman, where $L_\infty$ treats deformation. \\

$L_\infty$ algebra is a deformation of dgla, which is an analogue of Lie algebra for homological algebra
(consider Lie algebra consists of vector field, dgla consists of endomophism of chain complex $d:V\to V$ satisfying $d^2=0$). \\


\begin{definition}{Hochschild Complex}{}
\end{definition}


\vspace{0.5cm}
\noindent\rule{\textwidth}{0.4pt}
\vspace{0.5cm}


Next, I'll consider $A_\infty$ algebra/category and its relationship with algebraic geometry.
The question is if ncfp dg-manifold of $A_\infty$ algebra/category can be identified with formal scheme,
but the problem is the missing data $d_Q$ used for $A_\infty$ coherence.
I'll consider derived algebraic geometry, where we use cdga $(A,d)$, consisting differential structure on a ring $A$.

\begin{definition}{Formal Scheme}{}
A formal scheme is a locally ringed space $(X,\mathcal{O}_X)$ that is locally isomorphic to formal spectrum $Spf(A)$
of a complete topological ring $A$. Assuming $A$ is $I$-adically complete for an ideal $I\subset A$ s.t. $A = \varprojlim_{n} A/I^n$.
Then as a set, $|Spf(A)| = \{ \frak{p} \subset A  | \text{ \frak{p} is open and contains $I^n$} for some $n$\}$. 
as a remark, prime ideal $\frak{p}$ is defined to be closed in Zariski topology, but from different point of view,
it could be open subset in $I$-adic topology, different from Zariski. \\

For a scheme $X$, formal subscheme $\hat{X_L}$ of $L\subset X$ is given by a limit of sheaf 
$\hat{X_L} = \varprojlim_{n} \mathcal{O}_X / \mathcal{I}_L^n$ where $\mathcal{I}_L$ is an ideal sheaf of $\mathcal{O}_X$ for $L$.
\end{definition}


\begin{definition}{Non-Commutative Formal Pointed DG-Manifold}{}
For coalgebra derivation $Q: T^c (V[1]) \to T^c (V[1])$ s.t. $Q^2=0$ of $A_\infty$ structure where $V$ is a vector space,
there is a correspondence to formal scheme $(V,Q) \leftrightarrow (Spf(A), d_Q)$, where $V^\vee$ is a dual vector space,
and $A = \hat{Sym(V^\vee)}$, and $d_Q: Spf(A) \to Spf(A)$ is a differential s.t. $d_Q (f) = Q(f)$ for every function $f\in A$.
\end{definition}



\begin{definition} {Derived Formal Scheme}{}
\end{definition}


\begin{definition} {}{}
\end{definition}


I'll consider its connection to Fukaya category.

In fact, ncfp dg-manifold of $A_\infty$ algebra/category can be identified with formal scheme, not exactly the same. 


\begin{definition} {a}{}
\end{definition}


\begin{definition}{Formal Moduli}{}
$\frak{M}_L$ \\

$\hat{M}_L = \varprojlim \mathcal{O}_X/ I_L^{n+1}$
\end{definition}

\begin{definition}{Fukaya Category}{}
\end{definition}


\begin{definition}{Derived Formal Scheme}{}
$Spf(A)$
\end{definition}


Now that it looks very connection with ordinary algebaric geometry. Consider .

\begin{question} {}{}
\begin{itemize}
\item Is there any intersection theory property of formal pointed dg-manifold?
\end{itemize}
\end{question}


\vspace{0.5cm}
\noindent\rule{\textwidth}{0.4pt}
\vspace{0.5cm}


Finally, the proof of homological mirror symmetry conjecture is discussed from multi-dimensional perspective, 
including counting Gromow-Witten invariant, or Balmer spectrum, while generalizing formal geometry 
of non-unital $A_\infty$-algebra to $A_\infty$-category. \\



\tableofcontents

\section{Formal Moduli Problem}



\section{Balmer Spectrum}

If $A_\infty$ category is symmetric monoidal category, its homotopy category is tensor triangulated category.

tensor triangulated category \\


Balmer spectrum \\


stable homotmopy theory \\

\section{GW Invariant}


\nocite{*}
\printbibliography


\end{document}

